% !TeX root = ../tfg.tex
% !TeX encoding = utf8


\iffalse
    ################# ASUNCIONES CONDITIONAL MEAN INDEPENDENCE
    1. Introducción a la inferencia causal, las asunciones de conditional mean independence,
    positivity y exchangeability o SUTVA, y las potential outcomes y el estimado de la diferencia de riesgos y del odds ratio (OR).
    positivity
    exchangeability o SUTVA
    potential outcomes
    estimado de la diferencia de riesgos y del odds ratio (OR)
\fi

% Intro a la inferencia causal
\chapter{Introducción a la inferencia causal}
\section{Motivación - RCT}
Para entender el razonamiento que lleva al desarrollo de los modelos causales, los compararemos
con los más ubicuos RCT.

Un RCT, por sus siglas en inglés de \textit{Randomized Controlled Trial} es un experimento en el que,
dentro de la muestra elegida, se asigna a los miembros a los grupos de control o de tratamiento de forma aleatoria.
Esta asignación aleatoria puede ser simple, siguiendo una distribución de Bernouilli, o más elaborada.
En cualquier caso, conocemos la distribución de probabilidad de la variable aleatoria que representa si una unidad es tratada o no.

En el límite de una muestra de gran tamaño, esto lleva a que la única diferencia entre los dos grupos es si son tratados o no,
incluso en el caso de que haya variables confusoras.
Esto nos permite eliminar el efecto del sesgo de selección y de las variables confusoras.

Es cierto que en experimentos clínicos usuales, la muestra no suele ser tan grande como para ignorar
otros factores confusores y covariantes, por los cuales se intenta \textit{ajustar} para reducir ese sesgo.

No siempre se pueden realizar estudios RCT por varios motivos, entre ellos la propia viabilidad, la ética o el coste
de realizarlos, por lo que uno se ve obligado a hacer estudios observacionales.

En un estudio observacional, el investigador no puede controlar a qué unidades, pacientes o comunidades se le aplica el tratamiento,
por lo que se debe tener en cuenta la distribución del tratamiento en la población.

\begin{observacion}
    De una forma más elocuente a la presentada aquí introduce Rubin este problema en
    \cite[sección 2, párr. 6]{rubin2005}:

    \begin{quotation}
        One cannot compare medical compare medical costs for male
        smokers with medical costs for female non smokers and claim
        that this comparison is a causal effect of smoking without
        making heroic and unwarranted assumptions.
    \end{quotation}
\end{observacion}


\section{El experimento ideal}
El llamado experimento ideal \cite[sección 1.1]{targetedLearning} consiste en considerar una población y, en condiciones idénticas,
usar la población tanto como grupo de control y como grupo de tratamiento al mismo tiempo.
Comparar estos dos grupos nos permitirá hacer conclusiones causales. Dejamos la causalidad como un término
no definido de momento.
% TODO: definicion causalidad?

Este experimento es, naturalmente, imposible, pues son situaciones incompatibles. Este es el "problema fundamental de la inferencia
causal" \cite[sección 2.1]{mci}

% Con asunciones adicionales, construiremos un modelo causal a partir de un modelo estadístico.


\section{Estimandos de interés}
Vamos a considerar la Diferencia de Riesgo (RD, por sus siglas en inglés), el
Riesgo Relativo (RR) y el Odds Ratio (OR). % TODO: translation OR

En las siguientes definiciones consideramos una distribución $P$ de una variable bidimensional
$(A, Y)$, donde tanto $A$ como $Y$ son variables binarias. $A$ es un factor de riesgo
e $Y$ es la variable de respuesta.

\begin{definicion}[Diferencia de Riesgo] \label{def:riskdiff} % TODO: source
    Definimos la diferencia de riesgo (RD) como:
    $$
        RD = P(Y = 1 | A = 1) - P(Y = 1 | A = 0)
    $$
\end{definicion}

\begin{definicion}[Riesgo Relativo] \label{def:relrisk} % TODO: source
    Definimos el Riesgo Relativo (RR) como:
    $$
        RR(P) = \frac{P(Y = 1 | A = 1)}{P(Y = 1 | A = 0)}
    $$
\end{definicion}

\begin{definicion}[Odds Ratio] \label{def:oddsratio} % TODO: source
    Definimos el Odds Ratio (OR) como:
    % OR = RR * \dfrac{P(\neg E|C)}{P(\neg E | I)} = RR * \dfrac{1 - P(E|C)}{1 - P(E|I)}
    \begin{align*}
        OR & = \frac{\frac{P(Y = 1 | A = 1)}{P(Y = 0 | A = 1)}}{\frac{P(Y = 1 | A = 0)}{P(Y = 0 | A = 0)}} \\
        OR & = \frac{P(Y = 1 | A = 1)}{P(Y = 1 | A = 0)} * \frac{P(Y = 0 | A = 0)}{P(Y = 0 | A = 1)}
    \end{align*}
\end{definicion}

\begin{corolario}[Relación entre OR y RR]
    $$
        OR = RR * \dfrac{P(Y=0|A=0)}{P(Y=0|A=1)}
    $$
\end{corolario}

\begin{corolario}[Caso de enfermedad rara]
    Si $P(Y = 1 | A = a) \approx 0 \; \forall a$, entonces $OR \approx RR$.
\end{corolario}

Estudiaremos los estimadores asociados en \autoref{sec:estimadores_asociados}
